Scrum Framework
The project was developed using the Scrum framework to support iterative development, teamwork, and continuous progress throughout the implementation of the heat production optimization system. The work was divided into six sprints, with a different Scrum Master assigned for each sprint. This rotation was a team decision made to ensure that every member gained experience with sprint coordination and agile project management.
Scrum Roles
The Scrum Master was responsible for coordinating meetings, preparing agendas for supervisor meetings, writing meeting logs, and ensuring that Scrum practices were followed throughout each sprint. This included monitoring ticket progress, identifying blocked tasks, ensuring sprint goals were achieved, and making sure improvements from retrospectives were implemented. The development team collectively handled implementation tasks, sprint planning, and estimation. Two supervisors supported the project: a main supervisor and a substitute supervisor. The main supervisor provided weekly feedback, while the substitute supervisor assisted when needed.
Meetings and Ceremonies
The team followed standard Scrum ceremonies, including Sprint Planning, Daily Stand-ups, Sprint Reviews, and Sprint Retrospectives.
Sprint Planning included task breakdown and estimation using Planning Poker, which allowed the team to collaboratively evaluate task complexity and workload.
Sprint Reviews were held to present completed work and receive feedback from the supervisor.
Sprint Retrospectives were used to identify improvements for the next sprint.
Internal Scrum meetings were held every Monday and Wednesday, typically on-site in Løkken. These meetings included daily stand-ups where team members reported progress, discussed blockers, and assigned new tasks. When in-person meetings were not possible, communication and coordination were handled via Discord.
Weekly meetings with the main supervisor were held every Friday, either online or in person, to present progress, discuss technical challenges, and receive feedback.
Tools Supporting Scrum
Jira was used as the main project management tool to manage the product backlog, track sprint progress, and organize tasks. Jira was chosen instead of Trello due to its stronger support for task dependencies, sprint tracking, and structured workflow management in larger projects.
Excalidraw was used to visually map system architecture, workflows, and planning decisions, improving shared understanding within the team.
Adjustments to Scrum
Although the project followed core Scrum principles, several adjustments were made to fit the academic environment and team constraints:
Sprint length and workload were adapted to match academic deadlines and exam periods.
The Scrum Master role rotated between team members each sprint, instead of being a fixed role.
Some Scrum ceremonies were combined or shortened when schedule conflicts occurred.
Communication tools (Discord) were used in addition to in-person meetings to support remote collaboration when necessary.
These adjustments ensured that Scrum remained practical and effective within the constraints of a student development project. 
